\section{Discussion}
\label{sec:experiments-discussion}

% =====================================================================
%   STUB. Sub-section structure follows OUTLINE_TODO §5.5. Each
%   subsection is currently a heading + a brief note describing the
%   material that will land there, drafted from TODO.md where prose
%   already exists.
% =====================================================================

\subsection{The $\phi$-saturation pathology}
\label{subsec:discussion-phi-saturation}

% Material draft (see TODO.md, "Discussion" entry):
%
% The terminal Softmax layer in the M3N LeNet-5 variants
% (Section ref{subsec:method-architecture-backbones}) bounds the
% unary potentials so that the per-example dual variables phi^i
% cannot inflate to compensate for an over-confident backbone.
% Without the Softmax, training empirically degenerates: phi^i grows
% unboundedly, the pairwise term saturates against it, and the
% predictor reduces to a per-cell classifier with no constraint
% propagation. Figure: phi_norm trajectory from
% history.diagnostics across runs with/without the Softmax tail.

\subsection{$\mathrm{lr\_pairwise}$ rebalancing: CNN-vs-$W$ training
asymmetry at low $N$}
\label{subsec:discussion-lr-pairwise}

% Material draft (preserved in TODO.md under "lr_pairwise rebalancing"):
%
% Mechanism (durable). At low N the CNN converges much faster than W
% (CNN has ~60k params + ~36 clue images per puzzle; W has 81 params
% and needs the constraint structure to emerge from cross-batch
% gradient pressure). Default identical learning rates leave W
% under-trained --- the "0.35 plateau" is the structural signature
% of "CNN done, W partially-learned" (clue cells correct, blanks
% ~37% accurate via soft-constraint propagation). A larger
% lr_pairwise rebalances by giving W more gradient juice while
% leaving the CNN trajectory unchanged.
%
% Symptom (reproducible). With lr_pairwise = lr_model = 0.001 (visual
% Sudoku, N=500, lambda=1e-4, B=32): training plateaus at hamming
% approx 0.35 by epoch 100 and crawls. Final hamming after 1000
% epochs approx 0.20, test 0/1 approx 0.83.
%
% Effect (validate post-§5.3 rerun). With lr_pairwise = 0.01 (10x
% lr_model): plateau is crossed by epoch 10. By epoch 25
% hamming = 0.051, test 0/1 = 0.43. Pre-rerun numbers --- confirm
% or revise once notebook 05 finishes with the loader fix.

\subsection{Failure modes}
\label{subsec:discussion-failure-modes}

% Three concrete cases, all populated post-§5.3:
%   - OCR-baseline puzzles where solver returns infeasible:
%     solver_feasibility metric, frequency vs OCR error rate
%   - M3N predictions that violate Sudoku row/column/box constraints:
%     rate, where they cluster
%   - Per-puzzle correlation between blank-cell density and 0/1 error

\subsection{Limitations and threats to validity}
\label{subsec:discussion-limitations}

% N range, single dataset (Kaggle bryanpark/sudoku), BP vs exact
% decoder gap, finite-seed variance.

\subsection{Connections to broader literature}
\label{subsec:discussion-connections}

% Franc-Yermakov 2021, Kadlec 2022, SATNet (Wang et al. 2019).
